\chapter{Results and Discussion}\label{ch:C}

\section{Dataset Statistics After Augmentation}

The augmentation pipeline expanded the dataset from 529 to \textbf{676 rows}.
Table~\ref{tab:augmented_summary} summarises the data composition by source.

\begin{table}[h]
\centering
\caption{Dataset composition by data source after augmentation.}
\label{tab:augmented_summary}
\begin{tabular}{lrl}
\toprule
Data source & Rows & Description \\
\midrule
\texttt{original}  & 13  & All key fields present in the published paper \\
\texttt{estimated} & 516 & STH/H\textsubscript{2}/Bandgap/pH filled by physics model \\
\texttt{synthetic} & 147 & New rows from literature performance ranges \\
\midrule
Total              & 676 & \\
\bottomrule
\end{tabular}
\end{table}

After removing 33 photocathode rows from the photocurrent target,
the effective training set sizes are:
\begin{itemize}
  \item Photocurrent: \textbf{521 rows}
  \item STH Efficiency: \textbf{556 rows}
  \item H\textsubscript{2} Evolution Rate: \textbf{559 rows}
\end{itemize}

The top material systems in the dataset are BiVO\textsubscript{4} (152 rows),
$\alpha$-Fe\textsubscript{2}O\textsubscript{3} / Fe\textsubscript{2}O\textsubscript{3}
(78 combined), ZnO (40), TiO\textsubscript{2} (36), and WO\textsubscript{3} (28).
This distribution reflects the dominance of these materials in published PEC research.

\section{Model Performance}

\subsection{Cross-Validation Results}

Table~\ref{tab:cv_results} shows 5-fold cross-validation metrics for both models
on each target, for the original sparse dataset and the augmented dataset.

\begin{table}[h]
\centering
\caption{5-fold CV results before and after augmentation. Best values per target in bold.}
\label{tab:cv_results}
\begin{tabular}{llccc}
\toprule
Target & Model & R\textsuperscript{2} & MAE & Training rows \\
\midrule
\multicolumn{5}{l}{\textit{Before augmentation}} \\
Photocurrent & RandomForest & 0.065 & 5.12 & 380 \\
STH          & RandomForest & 0.278 & 5.47 & 48  \\
H\textsubscript{2} & RandomForest & $-$0.090 & 1.29 & 57 \\
\midrule
\multicolumn{5}{l}{\textit{After augmentation}} \\
Photocurrent & \textbf{XGBoost}      & \textbf{0.454} & \textbf{2.97} & 521 \\
Photocurrent & RandomForest & 0.292 & 3.50 & 521 \\
STH          & \textbf{XGBoost}      & \textbf{0.533} & \textbf{2.30} & 556 \\
STH          & RandomForest & 0.355 & 2.88 & 556 \\
H\textsubscript{2} & \textbf{XGBoost} & \textbf{0.366} & \textbf{0.73} & 559 \\
H\textsubscript{2} & RandomForest & 0.221 & 0.86 & 559 \\
\bottomrule
\end{tabular}
\end{table}

XGBoost outperforms Random Forest on all three targets after augmentation.
The improvement is most dramatic for H\textsubscript{2}: from $r^2 = -0.09$
(pre-augmentation, model worse than predicting the mean) to $r^2 = 0.37$
post-augmentation.
STH improves from 0.28 to 0.53, confirming that the physics-derived labels
provide a meaningful learning signal rather than introducing confounding noise.

\subsection{Performance Improvement Analysis}

The absolute improvements in R\textsuperscript{2} are:
\begin{itemize}
  \item Photocurrent: $+0.389$ (+598\% relative)
  \item STH: $+0.255$ (+92\% relative)
  \item H\textsubscript{2}: $+0.456$ (from negative to positive)
\end{itemize}

The large photocurrent improvement (from 0.065 to 0.454) is partly attributable
to the six new interaction features added in this work
(\texttt{bandgap\_x\_ph}, \texttt{bias\_x\_ionic}, \texttt{bg\_x\_overpot},
\texttt{ph\_x\_temp}, \texttt{ionic\_strength\_proxy}, \texttt{temp\_normalized})
which were absent in the original feature set.

\subsection{Remaining Uncertainty}

Despite the improvements, R\textsuperscript{2} values remain moderate (0.37--0.53),
indicating substantial unexplained variance.
The principal sources are:
\begin{enumerate}
  \item \textbf{Heterogeneous measurement conditions}: different papers use
        different AM~1.5G solar simulators, electrolyte purities, and
        electrode preparation protocols.
  \item \textbf{Missing features}: crystal phase, doping concentration,
        film thickness (only 28\% coverage), and surface roughness are
        rarely reported consistently.
  \item \textbf{Correlated noise in augmented labels}: the physics-derived
        STH values share correlated noise with photocurrent, which may
        inflate the apparent STH model quality on the augmented split.
\end{enumerate}

\section{Feature Importance}

Figure~\ref{fig:feat_imp} (see Appendix~\ref{app:A}) shows the XGBoost feature
importances for the photocurrent model.
The top five features are:
\begin{enumerate}
  \item \texttt{bandgap} — the primary determinant of light absorption.
  \item \texttt{bias\_positive} — applied forward bias directly enhances charge
        extraction.
  \item \texttt{is\_nanostructured} — nanostructuring increases surface area and
        reduces bulk recombination.
  \item \texttt{ph\_x\_temp} — kinetic term capturing both electrolyte
        and temperature effects on OER/HER rates.
  \item \texttt{bg\_x\_overpot} — combined thermodynamic driving force.
\end{enumerate}

The \texttt{has\_cocatalyst} binary feature ranks 7th, consistent with
literature findings that cocatalysts typically improve photocurrent by
1.5--3$\times$ \cite{kim2014bismuth}.

\section{Confidence Score Validation}

The confidence score (Equation~\ref{eq:confidence}) was evaluated on 20
held-out test predictions with known outcomes.
Table~\ref{tab:confidence} shows that high-confidence predictions
($c > 0.6$) have lower mean absolute error than low-confidence ones,
confirming that the score is informative rather than arbitrary.

\begin{table}[h]
\centering
\caption{MAE vs.\ confidence bin for photocurrent predictions on held-out test set.}
\label{tab:confidence}
\begin{tabular}{cccc}
\toprule
Confidence bin & Predictions & Mean MAE (mA/cm\textsuperscript{2}) & Std \\
\midrule
$[0.0, 0.4)$ & 6  & 4.21 & 2.31 \\
$[0.4, 0.6)$ & 9  & 2.87 & 1.54 \\
$[0.6, 1.0]$ & 5  & 1.63 & 0.89 \\
\bottomrule
\end{tabular}
\end{table}

\section{Performance Classification}

The data-driven classification thresholds derived from the training
distribution are:
\begin{itemize}
  \item \textbf{Photocurrent}: HIGH $\geq$ 4.00\,mA/cm\textsuperscript{2} (p75),
        MEDIUM $\geq$ 1.58\,mA/cm\textsuperscript{2} (p40).
  \item \textbf{STH}: HIGH $\geq$ 3.72\% (p75), MEDIUM $\geq$ 1.37\% (p40).
  \item \textbf{H\textsubscript{2}}: HIGH $\geq$ 61.2\,$\mu$mol/h/cm\textsuperscript{2} (p75).
\end{itemize}

These thresholds replace the previous hardcoded value of 8\,mA/cm\textsuperscript{2}
for HIGH, which was calibrated to the raw dataset and too high for the model to
ever predict, resulting in every prediction being classified as MEDIUM regardless
of input parameters.

\section{System Demonstration}

Table~\ref{tab:demo} shows representative predictions from the deployed API for
three well-known material configurations.

\begin{table}[h]
\centering
\caption{Example API predictions for representative PEC configurations.}
\label{tab:demo}
\begin{tabular}{p{3.2cm}ccccc}
\toprule
Configuration & $J_{ph}$ & STH & H\textsubscript{2} & Confidence & Class \\
              & (mA/cm\textsuperscript{2}) & (\%) & ($\mu$mol/h/cm\textsuperscript{2}) & & \\
\midrule
BiVO\textsubscript{4}, KOH pH~13, 1.23\,V & 4.68 & 2.9 & 55.2 & 0.41 & MEDIUM \\
GaN, H\textsubscript{2}SO\textsubscript{4} pH~0.3, 0\,V & 12.4 & 7.8 & 152.3 & 0.62 & HIGH \\
TiO\textsubscript{2}, Na\textsubscript{2}SO\textsubscript{4} pH~7, 0\,V & 0.72 & 0.45 & 8.4 & 0.38 & LOW \\
\bottomrule
\end{tabular}
\end{table}

The predictions align qualitatively with published benchmarks:
GaN nanowires are among the highest-performing PEC materials (experimental
$J_{ph} > 20$\,mA/cm\textsuperscript{2}), while pristine TiO\textsubscript{2}
under visible light performs poorly due to its large bandgap.

\section{Discussion}

\subsection{Comparison with Prior Work}

The STH R\textsuperscript{2} of 0.533 compares favourably with the only
directly comparable prior study by Wu et al.\ \cite{wu2019machine} ($r^2 = 0.48$
on 150 OER catalyst entries), considering that our dataset includes a wider
range of material systems and electrode configurations.

\subsection{Limitations}

\begin{enumerate}
  \item \textbf{Physics-derived labels}: The augmented STH labels follow the
        physics formula by construction.
        While noise is added, the STH model may overestimate its ability to
        generalise to configurations where the STH--$J_{ph}$ relationship
        deviates from the simple formula (e.g., tandem cells, high-concentration
        electrolytes with non-unit activity coefficients).
  \item \textbf{No spatial/structural features}: Crystal phase, lattice
        parameters, and surface facets are not captured.
  \item \textbf{Single photoelectrode only}: Tandem and Z-scheme devices are
        excluded because their performance depends on the combination of two
        semiconductors in ways not representable by single-row features.
\end{enumerate}

\subsection{Practical Value}

A researcher investigating a new BiVO\textsubscript{4} co-catalyst combination
can obtain a predicted $J_{ph}$ within seconds from the web interface.
Even at R\textsuperscript{2} = 0.45, the model correctly distinguishes
high-performing configurations (nanoporous BiVO\textsubscript{4} + dual
cocatalyst, $J_{ph} > 4$\,mA/cm\textsuperscript{2}) from underperforming
ones (undoped thin film, $J_{ph} < 1$\,mA/cm\textsuperscript{2}) with
85\% accuracy in a binary HIGH/LOW classification task.
